\section{The two-ended double-null interior}\label{sec:interior}

The exterior estimates supply data on the two late ends of a fixed,
deep spacelike cylinder. Its bounded middle is supplied by finite-time
evolution of the bridge. We construct the future double-null development
of these cylinder data, with uniform low-order control and arbitrarily
small exponential losses at every fixed higher order. The latter bounds
will be used on finite regions whose optical length tends to infinity.
The double-null framework follows the interior-stability method of
Dafermos--Luk \cite{DafermosLuk2017}; the entry and wedge estimates
required here are proved below. Throughout, the physical Kerr parameters
lie in a sufficiently small compact neighborhood of the fixed pair
$0<\mathfrak a<M$.

The null estimates below require exact profiles with one common
exponential rate and summable entry errors. We establish these inputs
before deriving the characteristic energy estimates. The profile depth will be chosen before the
inner computing boundaries, exterior coordinate constants, and data
smallness thresholds. No higher-order seminorm is required to be small.

We write a bound as $e^{o(T)}$ if, for every $\delta>0$, it is at
most $C_\delta e^{\delta T}$ for $T\ge1$. The constant may depend
on the fixed smooth datum, derivative order, and fixed domain
parameters. No uniformity in derivative order is asserted.

\subsection{Gauge and exact profiles}

Let $y_0,y_1$ be future-increasing optical functions and set
\begin{equation}\label{in:gauge}
\begin{gathered}
 s=\frac{y_0+y_1}{2},\qquad x=\frac{y_1-y_0}{2},\qquad
 \mathcal S=\{s=0\},\qquad q=q_0e^{-2\kappa s},\\
 g=-2a\,dy_0dy_1+
 \gamma_{AB}(d\theta^A-b^A dy_0)(d\theta^B-b^Bdy_0),
 \qquad a=qA.
\end{gathered}
\end{equation}
Here $\kappa>0$ is fixed, for example equal to the background inner
surface gravity in absolute value; $q_0>0$ will be chosen below;
$A>0$; $\gamma$ is a Riemannian metric on $\mathbb S^2$; and $b$ is
a vector field on $\mathbb S^2$. Angular coordinates are understood in
a finite smooth atlas. The symbols $a,A,q$ in this section denote the
null factor, its rescaling, and its exponential scale, respectively.
The rate $\kappa$ here is unrelated to the exterior radius-decay
exponent denoted by the same letter.

For distinct $i,j\in\{0,1\}$, write
\[
 B_0=b,\qquad B_1=0,\qquad
 L_i=\partial_{y_i}+B_i,\qquad
 \mathscr L_i=\partial_{y_i}+\mathcal L_{B_i}.
\]
The last operator acts on angular tensors. Define
\begin{equation}\label{in:variables}
\begin{gathered}
 \chi_i=\tfrac12\mathscr L_i\gamma,\qquad
 \tau_i=\operatorname{tr}_\gamma\chi_i,\qquad
 \ell_i=L_i\log a,\\
 \xi_0+\xi_1=d\log a,\qquad
 \xi_0-\xi_1=-a^{-1}\gamma\,\partial_{y_1}b,\\
 \beta_i=\operatorname{div}\chi_i-d\tau_i,\qquad
 c_i^{\mathrm{ang}}=\operatorname{curl}\xi_j,\qquad
 Q_i=(K,c_i^{\mathrm{ang}}),\qquad
 \mu_i=K-\operatorname{div}\xi_j.
\end{gathered}
\end{equation}
In this section $K$ is the Gaussian curvature of $\gamma$.
All angular differentials, contractions, divergences, and curls use
$\gamma$ unless explicitly indicated otherwise.

\begin{lemma}[Exact profiles with a common rate]\label{in:profiles}
For physical Kerr parameters in a sufficiently small compact
neighborhood of the fixed subextremal parameters, there are exact
profiles in the gauge \eqref{in:gauge}, depending smoothly only on
$(q,\theta)$ and the parameters and invariant under axial rotations,
such that $A,\gamma$ and their inverses extend smoothly to $q=0$,
$b=O(q)$, and
\begin{equation}\label{in:profile-orders}
 \chi_i,\ \ell_i+\kappa,\ \beta_i=O(q),\qquad
 \xi_i,\ K,\ c_i^{\mathrm{ang}},\ \mu_i=O(1).
\end{equation}
These estimates hold after every fixed angular or physical-parameter
derivative. Their constants can be fixed before $q_0$ is chosen.
The same statement holds with either optical direction chosen to have
zero angular shift.
\end{lemma}

\begin{proof}
Align the spin axes smoothly in the central trapped
block. Let $T_K$ be signed Boyer--Lindquist Killing time, and let
$(\vartheta,\varphi)$ denote the corresponding sphere variables.
For the current exact physical parameters, put
\[
 \Sigma_K=r^2+\mathfrak a^2\cos^2\vartheta,\quad
 \Delta=(r-r_-)(r-r_+),\quad
 r_*'=\frac{r^2+\mathfrak a^2}{\Delta},\quad
 \kappa_-=\frac{r_+-r_-}{2(r_-^2+\mathfrak a^2)}.
\]
The Boyer--Lindquist metric is
\begin{equation}\label{in:BL-metric}
\begin{split}
 g_K={}&-\frac{\Delta}{\Sigma_K}
        (dT_K-\mathfrak a\sin^2\vartheta\,d\varphi)^2
       +\frac{\Sigma_K}{\Delta}\,dr^2+\Sigma_K\,d\vartheta^2\\
 &+\frac{\sin^2\vartheta}{\Sigma_K}
        ((r^2+\mathfrak a^2)d\varphi-\mathfrak a\,dT_K)^2.
\end{split}
\end{equation}
Temporal and axial conventions may be reversed simultaneously when
passing between the two ends.

Solve the axisymmetric Hamilton--Jacobi equation
\begin{equation}\label{in:HJ}
 2(r^2+\mathfrak a^2)F_r+\Delta F_r^2
       +|d_\omega F|_{\mathrm{rd}}^2
       +\mathfrak a^2\sin^2\vartheta=0,
 \qquad F(r_-)=0,
\end{equation}
where the norm is that of the round sphere. The finite root is
noncharacteristic at the initial sphere, since its coefficient of
$F_r$ is $2(r_-^2+\mathfrak a^2)>0$. Smooth Hamilton--Jacobi evolution
therefore gives a parameter-dependent solution on a short radial
interval, uniformly over the chosen compact parameter set. The equation
is smooth on the whole sphere, and uniqueness preserves axisymmetry.
The inverse of \eqref{in:BL-metric} gives
\[
\begin{split}
 \Sigma_K g_K^{-1}
  \bigl(d(r_*+F\pm T_K),d(r_*+F\pm T_K)\bigr)
 ={}&\Delta\left(\frac{r^2+\mathfrak a^2}{\Delta}+F_r\right)^2
       +|d_\omega F|_{\mathrm{rd}}^2\\
 &-\frac{(r^2+\mathfrak a^2)^2}{\Delta}
       +\mathfrak a^2\sin^2\vartheta.
\end{split}
\]
Thus \eqref{in:HJ} is precisely the null equation.

Define
\begin{equation}\label{in:exact-coordinates}
 Q=e^{-2\kappa_-(r_*+F)},\qquad
 s=-\frac{\log(Q/q_0)}{2\kappa},\qquad
 \lambda_- =\frac{\kappa_-}{\kappa},\qquad
 x=\lambda_-T_K.
\end{equation}
Both $s-x$ and $s+x$ increase to the future near the inner horizon.
Indeed their sum has temporal gradient, with $s$ increasing as $r$
decreases. The parameter-smooth normalization
\[
 r_*=-\frac{1}{2\kappa_-}\log(r-r_-)+\text{smooth}
\]
shows that $Q$ is $(r-r_-)$ times a smooth strictly positive factor.
The factor $\lambda_-$ in both optical directions is what makes
$\kappa$ common to every profile.

Corotate the sphere by $-\Omega_-T_K$, where
$\Omega_- =\mathfrak a/(r_-^2+\mathfrak a^2)$.
In the resulting $(s,x,\omega)$ coordinates, coefficients are independent
of $x$ and invariant under axial rotations. To see their regularity,
the pairings of $\partial_{T_K}+\Omega_-\partial_\varphi$ with itself
and with angular vectors are $O(Q)$ by \eqref{in:BL-metric}.
The simple pole of the radial coefficient is canceled by
$r_s,\partial_\vartheta r|_s=O(Q)$, and $r_s/Q$ is smoothly
nonzero. Hence the angular cross coefficients with $ds,dx$ are $O(Q)$,
while the angular metric extends smoothly and positively to $Q=0$.
The absolute spacetime volume divided by $Q$ is smooth and positive,
as follows from the Boyer--Lindquist density $\Sigma_K$ relative to
round sphere measure. These are smooth sphere-tensor statements,
including at the poles; here $r_->0$ is essential.

The two-shift null decomposition consequently has shifts $O(Q)$ and
positive null factor $Q$ times a smooth positive function. Eliminate
the second shift by angular labels
$H(Q,\omega)$, with $H(0,\cdot)=\operatorname{Id}$.
After division by $Q$, their transport equation has the form
$-\kappa\partial_Q$ plus a smooth sphere vector field.
It yields axis-equivariant sphere diffeomorphisms and the one-shift
gauge \eqref{in:gauge}, with $q=Q$. Since $L_iq=-\kappa q$,
the definitions \eqref{in:variables} imply
\eqref{in:profile-orders}, with all asserted differentiated bounds.
The same construction eliminates the other shift after exchanging the
optical directions.
\end{proof}

Choose $q_0$ sufficiently small for the characteristic absorptions
below, uniformly for both angular-transport conventions. We also impose
the exact-profile depth condition used when returning the later packet
to the exterior. Since $\partial_\vartheta F=O(r-r_-)=O(q_0)$ at
entry and the spin stays bounded away from zero in the fixed
parameter neighborhood, we require the uniform strict inequality
\begin{equation}\label{in:profile-depth}
 \left(|\partial_\vartheta F|^2+
                   \mathfrak a^2\sin^2\vartheta\right)_{\mathrm{entry}}
 <4\mathfrak a^2.
\end{equation}
This supplies a strict positive margin for the radial potential
$(r^2+\mathfrak a^2)^2-\Delta
(|\partial_\vartheta F|^2+\mathfrak a^2\sin^2\vartheta)_{\mathrm{entry}}$
for $r\ge r_+$, since $\Delta<r^2$ and
$(r^2+\mathfrak a^2)^2/r^2\ge4\mathfrak a^2$.
Thus the central level $Q=q_0$ is
strictly between the horizons. The computing radius $r_b$ is then
chosen strictly deeper, with a margin before $r_-$; the exterior time
bends and compact constants may be fixed afterward. This order depends
only on exact physical profiles, not on smallness of high derivatives
of the eventual datum.
All fixed entry supports and their enlargements lie strictly between
the horizons, hence outside the physical exterior trapping neighborhood;
this separation persists after shrinking the parameter range.

\subsection{Moving coordinates and the entry estimates}

The estimates on the two exterior ends and a fixed finite evolution of
the middle bridge will supply the cylinder data. We derive their
trace bounds in prepared coordinates; Section~\ref{sec:comparison} will
identify these local evolutions geometrically and attach the resulting
cylinder to the complete bridge. Thus the entry estimate uses the
exterior bounds, while the later attachment uses geometric uniqueness.

The exterior stationary references are exact Kerr metrics in stationary
charts. On fixed spatial supports about the chosen entry depth, their
aligned Boyer--Lindquist variables have the form
\[
 T_K=\varsigma(\pi)t+P_K(\pi,y),
\]
with $r,\omega$ smooth in $(\pi,y)$. Here $\pi$ includes the required
exterior parameters, including the auxiliary stationary parameters
when necessary. The factor $\varsigma$ is nonzero and has opposite
signs on the two ends. Shrink the parameter range so that all these
supports remain within the fixed depth margins. Denote the associated
intrinsic physical parameters by $\mathfrak p(\pi)$.

\begin{lemma}[Preparation without growing phase errors]
\label{in:moving-coordinates}
For a time-dependent exterior parameter $\pi(t)$, define $s$ using
the instantaneous $Q$ in \eqref{in:exact-coordinates}, use the
instantaneous angular map $H$, and set
\begin{equation}\label{in:phases}
\begin{gathered}
 x=I(t)+\lambda_-P_K,\qquad
 \alpha_{\mathrm{rot}}=I_{\mathrm{rot}}(t)+\Omega_-P_K,\\
 I'=\lambda_-\varsigma,\qquad
 I_{\mathrm{rot}}'=\Omega_-\varsigma.
\end{gathered}
\end{equation}
The angular rotation is by $-\alpha_{\mathrm{rot}}$, after aligning
the axes. The Jacobian differs from the frozen Jacobian only by terms
containing positive parameter derivatives, with uniformly bounded
coefficients at fixed low order. No factor $t\pi'$ occurs.
Near $s=0$, the spheres of fixed $(s,x)$ are uniformly controlled
graphs over a stationary sphere, with graph times a bounded distance
from $I^{-1}(x)$ and uniformly nonsingular angular Jacobians.
\end{lemma}

\begin{proof}
This is the preparation in \Cref{in:phases}; we check its dependence on
the present exact family. At a time $t_0$, compare with the map frozen
at $\pi(t_0)$ and give its Killing-time and rotation products the
constant offsets
\[
 I(t_0)-(\lambda_-\varsigma)(\pi(t_0))t_0,
 \qquad
 I_{\mathrm{rot}}(t_0)-(\Omega_-\varsigma)(\pi(t_0))t_0.
\]
The actual and frozen maps then agree in value. The offsets change
neither the stationary nor the axisymmetric prepared profile.
Their differentiated discrepancies contain only positive parameter
jets with bounded smooth coefficients. Differentiating a phase gives
a bounded frequency or a parameter jet, never the product of an
elapsed time and a parameter derivative. Axis alignment precedes
rotation, so variations of the axis also have uniformly bounded
coefficients. The same comparison applies to the metric, with the
additional residual $g_{\mathrm{co}}-\overline G_{p,c}$.

For fixed $(s,x)$, first solve the $s$ equation for radius over a
stationary sphere. It is a uniformly small graph perturbation.
Along that level, the derivative of $x$ in time is
$I'+O(|\pi'|)$, bounded away from zero in absolute value.
The graph times therefore differ from $I^{-1}(x)$ by a bounded
amount. At fixed $x$, angular differentiation of the total rotating
phase costs $O(|\pi'|)$, by $dx=0$. The angular map is consequently
a uniformly nonsingular perturbation of an aligned rotation composed
with $H$. It is a sphere diffeomorphism; equivalently, finitely many
uniform graph charts give the same estimates.
\end{proof}

Fix the additive phase constants by agreement with the central exact
chart at zero perturbation. In bounded overlaps, join the auxiliary
coordinates by small smooth changes; small vector fields suffice for
the angular maps. The middle and these joins are contained in a fixed
finite-time bridge development. The resulting surface $s=0$ is
strictly spacelike, with the same uniform optical transversality.
Its precise attachment to the lower development uses ordinary
geometric Cauchy uniqueness and will be made in the comparison-domain
construction.

On each end define a physical parameter profile by
\begin{equation}\label{in:parameter-profile}
 \overline{\mathfrak p}(y_0)
   =\mathfrak p\bigl(\pi(I^{-1}(-y_0))\bigr),
\end{equation}
and join the two prescriptions smoothly, close to the background,
through a bounded interval. This is one angle-independent profile on
the entire wedge, depending only on $y_0$. A hat denotes the exact
Kerr value at $(q,\theta,\overline{\mathfrak p}(y_0))$,
including for a named differentiated geometric quantity. Thus a hat
means evaluation after differentiation at fixed physical parameters.
Differentiating the resulting hatted field adds only a defect in
direction $0$, bounded by the profile symbol constants times
\[
 \sigma_0(y_0)=|\overline{\mathfrak p}'(y_0)|,
 \qquad \sigma_1=0.
\]
This use of physical parameters keeps the profile constants independent
of the subsequent exterior chart and gauge choices.

\begin{proposition}[Cylinder data and their derivative budget]
\label{in:entry}
Start actual eikonal coordinates on the prepared cylinder with
$y_0=-x$, $y_1=x$, and transport its prepared angular labels along
$L_1$. For sufficiently small data in the common base neighborhood,
the sum of unit-$x$-cell supremum norms of the differences from frozen
profiles is at most an arbitrarily prescribed small number
$e_{\mathcal S}$ for
\begin{equation}\label{in:entry-orders}
\begin{array}{ll}
 A,\log A,\gamma,b,\chi_i,\ell_i,\xi_i
     &\text{in angular }H^6,\\
 \gamma&\text{in angular }H^7,\\
 \partial_x(\chi_i,\ell_i,\xi_i,b)|_{\mathcal S}
     &\text{in angular }H^2.
\end{array}
\end{equation}
Here $\partial_x$ is tangent to the cylinder at fixed prepared angle,
and the comparison in the last line uses frozen differentiation.
The sum of cell suprema of $\sigma_0$ is also at most
$C e_{\mathcal S}$. At each fixed higher order, these initial
quantities and their parameter jets on $|x|\le CT$ are bounded by
$e^{o(T)}$ for every fixed $C>0$.

The low smallness depends only on the common base-order neighborhood;
constants may depend on the already fixed $q_0$ and coordinate choices.
Higher-order constants may depend on the individual smooth datum.
\end{proposition}

\begin{proof}
We derive the entry trace from the mixed exterior estimates.
On the fixed graph supports, Proposition~\ref{co2:sliced} controls
mixed derivatives of
$g_{\mathrm{co}}-\overline G_{p(t),c(t)}$ through order eight,
and one further time derivative of that order, in local spacetime
$L^2$ with squared time weight $\langle t\rangle^{A_S}$,
where $A_S>1$. These supports are outside physical trapping.
Positive parameter derivatives through order eight have enlarged-cell
supremum bounds supplied by their through-order-nine weighted $L^2$
controls, again with exponent greater than one. The one-dimensional
$H^1$ inequality and weighted Cauchy--Schwarz make their cell sums
small. Bounded graph-time offsets cost the corresponding local
parameter variations.

The prepared frozen fields are the same smooth functions of intrinsic
parameters and $(q,\theta)$ at every time. After the value-matching
comparison of Lemma~\ref{in:moving-coordinates}, replacing the
instantaneous parameter by \eqref{in:parameter-profile} changes only
that parameter argument, with positive parameter jets charged when
differentiated. In particular it does not rotate a phase by
$t$ times a parameter variation.

For a residual derivative $f=D^{\alpha}
(g_{\mathrm{co}}-\overline G_{p,c})$, $|\alpha|\le7$, on an
enlarged fixed $(t,r)$ rectangle, the iterated one-dimensional Sobolev
inequality gives
\begin{equation}\label{in:trace}
 \int_{\mathbb S^2}\sup_{t,r}|f|^2\,d\omega
 \le C\sum_{\epsilon,\epsilon'\in\{0,1\}}
       \|\partial_t^\epsilon\partial_r^{\epsilon'}f
                    \|_{L^2(t,r,\omega)}^2.
\end{equation}
Even the mixed time-radial term lies in the stated budget: it uses
the full mixed order-eight estimate for the time derivative of the
metric residual, not a ninth-order pointwise metric trace.
The graph Jacobians and the bounded overlap of enlarged cells transfer
this estimate to the prepared spheres. If $a_n$ denotes the resulting
cell norm, then
\[
 \sum_n a_n
 \le\left(\sum_n\langle n\rangle^{-A_S}\right)^{1/2}
      \left(\sum_n\langle n\rangle^{A_S}a_n^2\right)^{1/2}.
\]
The first factor is finite, and the second is small by the sliced
estimate.

The first and second initial jets of the actual optical and angular
chart are smooth functions of the auxiliary metric through order one:
solve the noncharacteristic eikonal and angular-transport equations
for normal derivatives. At exact freezing these jets agree with the
auxiliary ones. Further initial jets require only ordinary
derivatives of these identities. The geometric fields in
\eqref{in:entry-orders} therefore use metric derivatives through
order seven: $A,\gamma,b$ use order zero and the one-derivative
geometric fields use order one before angular differentiation.
The coordinate transformations through this order use parameter jets
through order eight. Sphere Sobolev products, with spare angular
suprema on low residual factors, and \eqref{in:trace} give exactly
the displayed norms. Strict optical transversality requires only low
suprema after $q_0$ and the coordinate constants have been fixed.

In the bounded middle and enlarged joins, finite-time spacetime
estimates through order nine supply the same trace. Use common
identified coordinates for the lower computations; any order loss in
the identifying map is assigned to the pulled-back metric bounds in
that finite-time matching. Smooth fixed coordinate prescriptions and
small parameter-dependent changes suffice to interpolate the
preparation there. At higher orders the identical argument uses the
subexponential exterior estimates on expanding intervals and finite
smooth constants in the middle. This proves every assertion.
\end{proof}

\subsection{The null equations and uniform low estimates}

The entry estimate supplies summable errors along the initial cylinder.
We now propagate those errors through the wedge using only the vacuum
equations in \eqref{in:gauge}, the exact profile orders, and the entry
bounds. The sums along each optical direction will control transport
coefficients on arbitrarily long null segments; their vanishing tails
will later give the small exponential rates at higher orders.

With curvature convention $R(X,Y)=[\nabla_X,\nabla_Y]-\nabla_{[X,Y]}$,
the vacuum identities in this gauge are
\begin{equation}\label{in:vacuum}
\begin{aligned}
 L_i\tau_i&=\ell_i\tau_i-|\chi_i|^2,\\
 \mathscr L_i\xi_j&=\beta_i+\tau_i\xi_i,\\
 \mathscr L_j\chi_i
   &=a\operatorname{Sym}\nabla\xi_j-\tfrac a2K\gamma
      +a\xi_j^{\otimes2}
      -\tfrac12(\tau_i\chi_j+\tau_j\chi_i)
      +\chi_i\chi_j+\chi_j\chi_i,\\
 L_j\ell_i
   &=aK-\chi_i\mathbin{\cdot}\chi_j+\tau_i\tau_j
      +a(2\xi_i\mathbin{\cdot}\xi_j-|\xi_i|^2).
\end{aligned}
\end{equation}
\paragraph{Derivation of the null equations.}
For angular vectors $E,H$, Koszul's formula and
$[L_i,L_j]=a(\xi_i-\xi_j)^\sharp$ give
\begin{equation}\label{in:connection-identities}
\begin{aligned}
 \nabla^g_{L_i}L_i&=\ell_iL_i,&
 \nabla^g_E L_i&=\chi_i(E)^\sharp+\xi_i(E)L_i,\\
 \nabla^g_{L_i}L_j&=a\xi_i^\sharp,&
 \nabla^g_EH&=\nabla_EH+
      a^{-1}\{\chi_i(E,H)L_j+\chi_j(E,H)L_i\}.
\end{aligned}
\end{equation}
The same-index null Ricci contraction is
$-L_i\tau_i+\ell_i\tau_i-|\chi_i|^2$, giving Raychaudhuri.
Direct calculation also gives
\[
 g(R(L_i,E)L_j,H)
 =(\mathscr L_i\chi_j-\chi_j\chi_i
       -a\nabla\xi_i-a\xi_i^{\otimes2})(E,H).
\]
The angular Gauss contraction is
\[
 K\gamma+a^{-1}
 \{\tau_i\chi_j+\tau_j\chi_i-\chi_i\chi_j-\chi_j\chi_i\}.
\]
Thus $\Ric_{EH}=0$ gives the sum of the two cross equations for
$\chi$.  Their difference follows from
\[
 \mathscr L_j\chi_i-\mathscr L_i\chi_j
 =\tfrac12\mathcal L_{[L_j,L_i]}\gamma
 =a\,\operatorname{Sym}\nabla(\xi_j-\xi_i)
     +a(\xi_j^{\otimes2}-\xi_i^{\otimes2}),
\]
where $da=a(\xi_i+\xi_j)$ was used.  Half the sum and difference give
the cross-shape equation of \eqref{in:vacuum}.

In $\Ric_{iE}$ the angular and null contractions are, respectively,
$\beta_i-\chi_i\xi_i+\tau_i\xi_i$ and
$-\mathscr L_i\xi_j+\chi_i\xi_i$.  This proves the torsion equation.
In $\Ric_{ji}$ the angular contraction, after use of the cross
equation, is $aK-\chi_i\cdot\chi_j+\tau_i\tau_j$, while the remaining
contraction is
$a(2\xi_i\cdot\xi_j-|\xi_i|^2)-L_j\ell_i$.  This proves the lapse
equation.

The corresponding differentiated angular systems have the factor structure
\begin{equation}\label{in:angular-system}
\begin{aligned}
 \mathscr L_j\beta_i
   &={}\tfrac a2(dK+Jd c_i^{\mathrm{ang}})
      +a(K*\xi+\xi*\nabla\xi+\xi^3)
      +\chi_i*\nabla\chi_j+\chi_j*\nabla\chi_i,\\
 L_i(Q_i,\mu_i)
   &={}((\operatorname{div}\beta_i,
                       \operatorname{curl}\beta_i),0)
      +\chi_i*(K,\nabla\xi)+\nabla\chi_i*\xi.
\end{aligned}
\end{equation}
Here $Jd$ is the rotated angular differential and $*$ denotes the
indicated smooth angular metric contractions; every differentiated
factor type is retained. Also
\begin{equation}\label{in:bracket-shift}
 [L_i,L_j]=a(\xi_i-\xi_j)^\sharp,\qquad
 \partial_{y_1}b=a\gamma^{-1}(\xi_1-\xi_0).
\end{equation}
To verify this, apply divergence minus trace differential to the cross
equation.  On $\operatorname{Sym}\nabla\xi_j$ the second-derivative
part is $\tfrac12Jd\,\operatorname{curl}\xi_j$; the curvature
commutators give $K*\xi$.  Differentiating $a$ uses
$da=a(\xi_i+\xi_j)$, and varying the angular connection uses
$\mathscr L_j\gamma=2\chi_j$.  These give exactly the displayed remaining
types.  Surface curvature variation gives
\begin{equation}\label{in:curvature-variation}
 L_iK=\operatorname{div}\beta_i-\tau_iK .
\end{equation}
For a one-form $\zeta$,
\[
 L_i\operatorname{div}\zeta
 =\operatorname{div}(\mathscr L_i\zeta)
   -2\chi_i^{AB}\nabla_A\zeta_B
   -(2\operatorname{div}\chi_i-d\tau_i)\cdot\zeta .
\]
Together with the torsion equation this cancels
$\operatorname{div}\beta_i$ in $L_i\mu_i$.  Variation of the area form
similarly gives
$L_i\operatorname{curl}\xi_j
 =\operatorname{curl}\beta_i+
   \operatorname{curl}(\tau_i\xi_i)-\tau_i\operatorname{curl}\xi_j$.
This proves \eqref{in:angular-system}.


Frozen hats satisfy these equations at fixed parameters. In differences,
parameter defects in direction $i$ contain only $\sigma_i$.

We first estimate the solution on finite past diamonds
\[
 \mathcal D_{C_0,C_1,s_*}
 =\{0\le s\le s_*,\ y_0\le C_0,\ y_1\le C_1\}
       \times\mathbb S^2,
 \qquad C_0+C_1>0,\quad s_*>0.
\]
At fixed $y_j$, the past coordinate segment ends at $y_i=-y_j$
on $\mathcal S$ and remains in this diamond. The upper null sides and
upper $s$ cuts are future boundaries for the energy currents. All norms
below are restricted to this domain. The slightly tilted enlargements
used for continuation will retain these outgoing flux signs.
Fix $0<\eta\ll\kappa$ and let
$w_{i,n}=e^{-\eta|y_i-n|}$, $n\in\mathbb Z$.
For nonnegative angular norm functions $Y,Z$, define
\begin{equation}\label{in:flux-norms}
\begin{aligned}
 [Y]_{i,n}
   &=\sup_{y_j}\|w_{i,n}Y\|_{L^2(dy_i)}
       +\sup_s\|w_{i,n}Y\|_{L^2(dx)},\\
 \{Z\}_{i,n}
   &=\sup_{y_i}\|w_{i,n}\sqrt q\,Z\|_{L^2(dy_j)}
       +\sup_s\|w_{i,n}\sqrt q\,Z\|_{L^2(dx)}
       +\|w_{i,n}\sqrt q\,Z\|_{L^2(dy_0dy_1)}.
\end{aligned}
\end{equation}
Every integral and supremum carries the domain restriction. Angular
$H^k$ norms in fixed controlled charts are written $|\cdot|_k$.
For actual-minus-hat differences, denoted by $\delta$, set
\begin{equation}\label{in:low-variables}
\begin{gathered}
 Y_i=|\delta\beta_i|_5,\qquad
 Z_i=|\delta(Q_i,\mu_i)|_5,\qquad
 R_i=|\delta\chi_i|_6,\qquad
 X_j=|\delta\xi_j|_6,\qquad d_i=|\delta B_i|_6,\\
 E=\sum_{i,n}([Y_i]_{i,n}+\{Z_i\}_{i,n}),\qquad
 R=\sum_{i,n}[R_i]_{i,n},\qquad
 X=\sum_{i,n}\{X_j\}_{i,n},\\
 D_b=\sum_{i,n}
        \|w_{i,n}\sup_{y_j}d_i\|_{L^2(dy_i)}.
\end{gathered}
\end{equation}
The indexing of $X_j$ in the $i$ norm is deliberate: $j=1-i$.
For the remaining low quantities put
\begin{equation}\label{in:low-point-budgets}
\begin{gathered}
 p_*=|\delta(\gamma,\log A)|_6+|\delta\xi|_5,\qquad
 P=\sup p_*+\sum_n\sup e^{-8\eta s}w_{0,n}p_*,\\
 g_i(y_i)=\sup_{y_j}
    \bigl(|\delta(\chi_i,\ell_i)|_4+|\delta B_i|_5\bigr),
 \qquad S=\sum_{i,n}\sup w_{i,n}g_i.
\end{gathered}
\end{equation}
The summations over cells are unsquared. They are needed both for
integrable transport coefficients and for their vanishing tails.
The paired angular equations estimate $E$. The remaining budgets
recover the shape and torsion fields ($R,X$), the shift ($D_b$),
the metric and lower torsion coefficients ($P$), and the integrable
directional envelopes ($S$) of $\delta\chi_i,\delta\ell_i,\delta B_i$.
These recoveries are coupled: the proof first reduces
their combined size to $E$ and the entry error, then closes the small
feedback into $E$.

\begin{proposition}[Uniform low characteristic estimates]
\label{in:low}
Choose $q_0$ sufficiently small in terms of the uniform profile
constants, and then choose $e_{\mathcal S}$ sufficiently small.
On every smooth finite past diamond satisfying the entry assumptions,
\begin{equation}\label{in:low-conclusion}
 E+R+X+D_b+P+S\le C(q_0)e_{\mathcal S}.
\end{equation}
The metrics $A,\gamma$ and their inverses remain uniformly bounded
and nondegenerate at the indicated low orders. The constants are
independent of the location and size of the diamond.
\end{proposition}

\begin{proof}
Bootstrap the budgets small and the metrics in a uniformly
nondegenerate neighborhood. We first derive the Hodge and transport
estimates for the auxiliary budgets. Trace damping supplies the
shape estimate needed in this group. We then estimate the paired
angular currents and perform the simultaneous absorption.
The calculation uses only the displayed profile and equation
hypotheses. The actual radiative factors at low orders have
size $C(q+g_i)$; nonradiative coefficients are bounded and their
differences cost $Cp_*$. An orientation change obeys
\begin{equation}\label{in:weight-switch}
 w_{i,n}\le e^{2\eta s}w_{j,-n}.
\end{equation}
Unit-cell convolution bounds the localized $L^2_x$ norm of $p_*$ by
$Ce^{10\eta s}$ times an $\ell^1$ envelope of total size $P$.
Positive powers of $q$ absorb this growth when $\eta$ is small.

The div--curl system for $\xi_j$ and the trace-free divergence
system for $\chi_i$ give
\begin{equation}\label{in:low-hodge}
 X\le C(E+\sqrt{q_0}P),\qquad
 R_i\le C\bigl(Y_i+|\delta(\tau_i/A)|_6
                         +|\delta\chi_i|_0+qp_*\bigr).
\end{equation}
Their constants use only the low angular ellipticity and coefficient
bounds. Coefficient differences acting on hats cost $p_*$, with the
additional $q$ in the shape estimate. Integrating the shift equation
in \eqref{in:bracket-shift}, with Cauchy--Schwarz and
\eqref{in:weight-switch}, yields
\begin{equation}\label{in:low-shift}
 D_b\le Ce_{\mathcal S}+Cq_0^c(X+P)
\end{equation}
for a fixed $c>0$. The cross transports for $\chi_i,\ell_i$ have
homogeneous coefficient $C(q+g_j)$ and remaining difference source
$Cq(p_*+g_j+\sigma_j)$. Their quadratic terms pair opposite
radiative orientations. The lower shift estimate has source $Cqp_*$.
Consequently
\begin{equation}\label{in:low-characteristic}
 S\le Ce_{\mathcal S}+Cq_0^c(P+S).
\end{equation}

The metric and torsion transports, using
$\partial_{y_1}\gamma=2\chi_1$,
$\partial_{y_1}\log A=\ell_1+\kappa$, and
$d\log A=\xi_0+\xi_1$, give
\begin{equation}\label{in:low-metric}
 P\le C(e_{\mathcal S}+E+R+D_b+S).
\end{equation}
The unweighted free sources integrate by the cell sums. For the
localized part, the target-to-source ratio of
$e^{-8\eta s}w_{0,n}$ is bounded along either past coordinate path;
the remaining weight change and the square root of the path length
are absorbed by this exponential factor, proving
\eqref{in:low-metric} with the current profile constants.

Raychaudhuri has the exact rescaled form
\begin{equation}\label{in:trace-equation}
 (L_i+\kappa)(\tau_i/A)=-A^{-1}|\chi_i|^2.
\end{equation}
The differentiated difference has source
$C((q+g_i)R_i+q^2p_*+q(\sigma_i+d_i))$,
and the tame transport commutators additionally cost
\[
 C(q+g_i)|\delta(\tau_i/A)|_6
                 +C d_i(g_i+qp_*).
\]
The first term is absorbed in the trace damping under the small
bootstrap. The second uses the small low envelopes and the shift
budget; its contribution is included in the
$\varepsilon_{\mathrm{boot}}(P+D_b+R)$ term below.
Exponential convolution
in $y_i$ preserves the localized flux norms: a displacement $h$
changes the weight by at most $e^{\eta h}$ and samples the earlier
$s-h/2$ slice. Combining this damping with
\eqref{in:low-hodge} gives
\begin{equation}\label{in:low-shape}
 R\le C(E+e_{\mathcal S}+S)
       +C(q_0^c+\varepsilon_{\mathrm{boot}})(P+D_b+R).
\end{equation}

It remains to close the paired angular energies. Commute
\eqref{in:angular-system} through angular order five and subtract
hats, retaining the actual principal operators. The source norms are
bounded by constants times
\begin{equation}\label{in:paired-sources}
\begin{aligned}
 G_i={}&q(X_0+X_1+Z_i+p_*+\sigma_j)
       +(q+g_j)(Y_i+R_i)+(q+g_i)(R_j+d_j),\\
 H_i={}&Y_i+R_i+d_i+\sigma_i+qp_*
       +(q+g_i)(X_0+X_1+Z_i).
\end{aligned}
\end{equation}
Indeed, angular principal commutators cost $qZ_i$ or $Y_i$;
coefficient differences on hats have the stated profile gain; and
unweighted shape products pair opposite orientations. Shift
commutators use the low radiative bounds. Parameter defects occur
only in their own direction and inherit the $q$ factor whenever the
hatted field is radiative. This verifies the source list \eqref{in:paired-sources}.

Use squared currents with weights $w_{i,n}^2$ on $Y_i$, flowing
along $j$, and $w_{i,n}^2a/2$ on $(Q_i,\mu_i)$, flowing along $i$.
Angular integration by parts cancels the principal pair. The latter
currents have damping because
\[
 L_i\log(aw_{i,n}^2)
     \le-\kappa+C(q+g_i)+2\eta.
\]
All future null and $s$ fluxes are positive. The squared flux is
bounded by the initial norm plus
\begin{equation}\label{in:paired-current}
 C\int w_{i,n}^2(Y_iG_i+qZ_iH_i)\,dy_i\,dy_j.
\end{equation}
Terms retaining a positive power of $q$ use slice
Cauchy--Schwarz and \eqref{in:weight-switch}. Same-orientation
$g$ terms use a null flux and $\int g_i$. The remaining
opposite-orientation terms obey \Cref{in:cross-flux}:
\begin{equation}\label{in:cross-flux}
\begin{aligned}
 \int w_{i,n}^2Y_i g_iR_j
  &\le C[Y_i]_{i,n}\|w_{i,n}g_i\|_{L^2(dy_i)}
                          \sum_m[R_j]_{j,m},\\
 \int q w_{i,n}^2Z_i g_iX_i
  &\le C\{Z_i\}_{i,n}\|w_{i,n}g_i\|_{L^2(dy_i)}
                          \sum_m\{X_i\}_{j,m}.
\end{aligned}
\end{equation}
For completeness, partition $y_j$ into unit cells. Two-dimensional
Cauchy--Schwarz uses the fixed-$y_j$ flux for the first energy factor
and the fixed-$y_i$ flux for the opposite factor; in the second line,
the spacetime $\sqrt q\,X_i$ norm on a cell is controlled by its
$j$-localized norm. Partial cells and past-diamond restrictions remain
inside the integrals. The same argument applies to $d_j$ in place of
$R_j$. Summing the unsquared budgets gives
\begin{equation}\label{in:low-energy}
 E\le Ce_{\mathcal S}
       +C(q_0^c+S)(E+R+D_b+P+X).
\end{equation}

To close the estimates, write
$\mathcal A_{\mathrm{aux}}=P+R+X+D_b+S$.
Substitute the Hodge bound for $X$, the shift and cross-transport
bounds for $D_b,S$, and the damped shape bound for $R$ into the
metric estimate for $P$. For some fixed $c'>0$, this gives
\[
 \mathcal A_{\mathrm{aux}}
 \le C(E+e_{\mathcal S})
      +C(q_0^{c'}+\varepsilon_{\mathrm{boot}})
                      \mathcal A_{\mathrm{aux}}.
\]
Here $c'$ may be decreased to include the $\sqrt{q_0}$ term in
\eqref{in:low-hodge}. Absorbing the last term yields
$\mathcal A_{\mathrm{aux}}\le C(E+e_{\mathcal S})$.
The paired-current estimate \eqref{in:low-energy}, using the
bootstrap bound for $S$, now becomes
\[
 E\le Ce_{\mathcal S}
       +C(q_0^{c'}+\varepsilon_{\mathrm{boot}})
                                      (E+e_{\mathcal S}).
\]
Choose $q_0$ first using only the profile and energy constants,
then the bootstrap size, to absorb this feedback. Finally choose
the entry size so that the resulting bounds strictly improve the
bootstrap, including nondegeneracy. This proves
\eqref{in:low-conclusion}.
\end{proof}

\subsection{The first pure longitudinal derivatives}

The low flux bounds control the cross derivatives in the null
equations. Normalized curvature also contains pure derivatives such
as $\mathscr L_i\chi_i$, which those equations do not directly bound
pointwise. Differentiating the cross equations gives transport
equations for these missing jets along the opposite direction, with
no additional small data norm.

\begin{proposition}[Pure jets and flank limits]\label{in:pure-jets}
Under the assumptions of Proposition~\ref{in:low}, the sum of
unit-$y_i$-cell suprema of
\begin{equation}\label{in:pure-budget}
 \sup_{y_j}\left(
 |\mathscr L_i\chi_i-\widehat{\mathscr L_i\chi_i}|_2
 +|L_i\ell_i-\widehat{L_i\ell_i}|_2\right)
\end{equation}
is at most $Ce_{\mathcal S}$. For $i=0$ the same holds for
$\partial_{y_0}b-\partial_{y_0}^{\mathrm{fr}}\widehat b$
in angular $H^2$. On a full wedge obtained by exhaustion,
$A^{\pm1},\gamma^{\pm1}$ are uniformly bounded in $H^6$ and
$\xi_i,a^{-1}\partial_{y_1}b$ in $H^5$, and
\begin{equation}\label{in:flank}
\begin{gathered}
 |b|_5+\sum_i|(\chi_i,\ell_i+\kappa)|_4\longrightarrow0,\\
 \sum_{k=0}^1\left(
 |\partial_{y_k}b|_2+
 \sum_i|\partial_{y_k}(\chi_i,\ell_i)|_2
 +|\partial_{y_k}\gamma|_4
 +|\partial_{y_k}\log a+\kappa|_4\right)\longrightarrow0
\end{gathered}
\end{equation}
as $s\to\infty$ and $|y_0|,|y_1|\to\infty$.
All displayed quantities are uniformly bounded elsewhere.
No smallness of the limiting angular metric or rescaled torsion is
asserted.
\end{proposition}

\begin{proof}
We use the pure transport equations.
Set $U_i=\mathscr L_i\chi_i$ and $V_i=L_i\ell_i$.
The identities
\begin{equation}\label{in:pure-identities}
\begin{gathered}
 \mathscr L_i\xi_j=\beta_i+\tau_i\xi_i,\qquad
 \mathscr L_i\xi_i=d\ell_i-\beta_i-\tau_i\xi_i,\\
 L_iK=\operatorname{div}\beta_i-\tau_iK,\qquad
 [\mathscr L_j,\mathscr L_i]
     =\mathcal L_{a(\xi_j-\xi_i)^\sharp}
\end{gathered}
\end{equation}
reduce derivatives of torsion, curvature, and angular connections to
already controlled fields. Differentiating the cross equations in
direction $i$ gives \Cref{in:pure-top}:
\begin{equation}\label{in:pure-top}
 \mathscr L_jU_i=-\tfrac12\tau_jU_i
                    +U_i\chi_j+\chi_jU_i+\mathcal F_i,
 \qquad
 L_jV_i=-U_i\mathbin{\cdot}\chi_j+\mathcal G_i.
\end{equation}
Neither remainder contains a new pure jet. Differentiating $a$
retains $a$; a differentiated angular connection uses $\chi_i$ and
its angular derivatives; and the bracket is an angular operator with
an $a$ factor. Differentiating unweighted shape products introduces
either this factor or a cubic product containing both orientations.
The strongest angular term is $\nabla\beta_i$, controlled in $H^2$
by the available $H^3$ bound for $\beta_i$.

After subtraction of frozen differentiated fields, the inhomogeneous
sources in $H^2$ cost
\[
 C\{q(p_*+g_i+g_j+\sigma_j)+g_i g_j\},
\]
and the homogeneous coefficient along $j$ is $C(q+g_j)$.
Only a top shape jet times $\chi_j$ recurs without a $q$ factor.
The initial pure jets are obtained from the tangential derivatives in
\eqref{in:entry-orders} and the known cross equations.
Localized Gronwall along $j$, with the same orientation and cell
estimates as above, proves \eqref{in:pure-budget}.

For the shift write $\zeta_\xi=\xi_1-\xi_0$ and
$C_b=a\gamma^{-1}\zeta_\xi=\partial_{y_1}b$.
The identity \Cref{in:pure-shift} is
\begin{equation}\label{in:pure-shift}
\begin{split}
 \partial_{y_1}\partial_{y_0}b
 ={}&a\{\ell_0\gamma^{-1}\zeta_\xi
       -2\gamma^{-1}\chi_0\gamma^{-1}\zeta_\xi
       +\gamma^{-1}(2\beta_0+2\tau_0\xi_0-d\ell_0)\}
       -[b,C_b].
\end{split}
\end{equation}
Every term retains $a$, and no new pure shift derivative appears on
the right. Its difference source costs $Cq(p_*+g_0)$ in $H^2$.
Since the profile depends only on $y_0$, differentiating its frozen
pure shift in direction $1$ adds no parameter defect. Integration
from its tangential/cross initial value proves the shift assertion.

Summable cell suprema vanish in the corresponding infinite tails.
Together with \eqref{in:profile-orders}, the point budget therefore
gives the first line of \eqref{in:flank}. The cross equations give
the opposite first derivatives, while \eqref{in:pure-budget} and
\eqref{in:pure-shift} give the pure ones. Converting Lie derivatives
to coordinate derivatives uses already controlled $b,d_\theta b$.
Finally,
\[
 \partial_{y_i}\gamma=2\chi_i-\mathcal L_{B_i}\gamma,
 \qquad
 \partial_{y_i}\log a+\kappa
       =\ell_i+\kappa-B_i\mathbin{\cdot}d\log a.
\]
These identities give the remaining $H^4$ limits. The uniform metric,
lapse, and torsion bounds follow from $P$ and smooth inversion, and
$a^{-1}\partial_{y_1}b=\gamma^{-1}(\xi_1-\xi_0)$ has the asserted
$H^5$ bound.
\end{proof}

\subsection{Higher estimates on finite and expanding regions}

We next distinguish two uses of higher regularity. On a fixed finite
diamond, arbitrary finite bounds suffice to restart the optical
development. On expanding diamonds, the packet construction requires
arbitrarily small exponential rates. The first estimate uses no
global tails; the second is applied only after the low estimates have
been exhausted to the full wedge.

\begin{proposition}[Higher estimates]\label{in:higher}
Assume the entry bounds of Proposition~\ref{in:entry} and the low
estimates on smooth finite past diamonds.
\begin{enumerate}
\item On every fixed finite diamond, all fixed ordinary coordinate
derivatives of $\gamma^{\pm1},\log A,b,a^{-1}\partial_{y_1}b$
have finite bounds depending only on the corresponding smooth initial
norms, the low bounds, and that diamond.
\item If the solution and its uniformly summed low estimates are
defined on the full wedge, fix $U\in\mathbb R$, $C>0$, and either
labeling $\{d,l\}=\{0,1\}$. On
\begin{equation}\label{in:high-domain}
 s\ge0,\qquad y_d\le U,\qquad y_l\le CT,
\end{equation}
every fixed ordinary coordinate derivative of
$\gamma^{\pm1},\log A,b,a^{-1}\partial_{y_1}b$ is bounded by
$e^{o(T)}$. The same bounds hold on all restrictions and coordinate-path
subsegments of this domain.
\end{enumerate}
Neither conclusion requires any higher-order initial norm to be small.
\end{proposition}

\begin{proof}
We derive high-order currents and recover the unweighted derivatives
from their directional estimates. All high variables below are actual fields,
not differences from profiles. Put
\[
 h_i=C_0(q+g_i),\qquad \rho_q=\sqrt q,
\]
where $C_0$ bounds all required low radiative norms. The functions
$h_i$ are bounded and have uniformly bounded directional integrals.
On a full wedge their tails become arbitrarily small when
$s,|y_i|\to\infty$.

Fix an angular order $m\ge5$. In the rest of the proof redefine
\begin{equation}\label{in:high-variables}
\begin{gathered}
 Y_i=|\beta_i|_m,\qquad
 Z_i=\sqrt q\,|(Q_i,\mu_i)|_m,\qquad
 R_i=|\chi_i|_{m+1},\qquad z_i=|\tau_i/A|_{m+1},\\
 d_0^*=D=|b|_{m+1},\qquad d_1^*=0,\qquad
 P_\gamma=1+|\gamma|_{m+1},\qquad
 P_i=|\xi_j|_m,\qquad P_*=P_\gamma+P_0+P_1.
\end{gathered}
\end{equation}
In particular, $Z_i$ already includes the lapse weight.
The quantity $P_*$ controls $\gamma^{-1}$ and $A^{\pm1}$ at the
corresponding angular orders, using the low bounds and
$d\log A=\xi_0+\xi_1$.

\paragraph{Angular currents.}
The tame sphere Hodge estimates give
\begin{equation}\label{in:high-hodge}
 R_i\le C_m(Y_i+z_i+h_iP_*),\qquad
 |\xi_j|_{m+1}\le C_m(Z_i/\sqrt q+P_*).
\end{equation}
These constants do not contain high coefficient norms. Indeed,
identify trace-free tensors using smooth $\gamma$-dependent bases,
freeze coefficients in sufficiently small charts, and commute the
first-order angular system. A highest coefficient placement costs
$C_m|\gamma|_{m+1}\|f\|_{C^1}$; interpolation absorbs the
penultimate field norm with a constant depending on $m$ and the low
bounds. In the shape estimate its low factor is $O(h_i)$.

Use the squared paired currents $Y_i^2$ in direction $j$ and
$Z_i^2$ in direction $i$, keeping their multipliers equal.
Add $z_i^2$ and $P_i^2$ in direction $i$, and $D^2,P_\gamma^2$
in direction $1$. Apart from favorable damping, their absolute
interaction costs are bounded by the following list:
\begin{equation}\label{in:high-costs}
\begin{aligned}
 C_mY_i\big[&\rho_q(Z_0+Z_1+P_*)
       +h_j(Y_i+R_i)+h_i(R_j+d_j^*)+h_i h_jP_*\big],\\
 C_mZ_i\big[&\rho_q(Y_i+R_i+d_i^*+P_*)
                                      +h_i(Z_0+Z_1)\big],\\
 C_mz_i h_i(&z_i+R_i+d_i^*+h_iP_*),\\
 C_mD\big[&\rho_q(Z_0+Z_1+\rho_qP_*)+h_1D\big],\\
 C_mP_\gamma(&R_1+h_1P_\gamma),\\
 C_mP_i(&Y_i+R_i+d_i^*+h_iP_*).
\end{aligned}
\end{equation}
We derive these costs from the actual equations already displayed.
Terms with $a$ retain
$\rho_q$ after pairing, by \eqref{in:high-hodge}; high angular
connection placements in $\chi*\nabla\chi$ cost $h_i h_jP_*$.
Shift commutators on $\beta_i$ cost $h_jY_i+h_id_j^*$, while
those on the weighted fields retain $\rho_q$. Metric and area
derivatives have coefficient $h_j$ on $Y_i^2$ and $h_i$ on
$Z_i^2$. The identity $L_i\log a=-\kappa+O(h_i)$ supplies the
favorable lapse damping. Trace and shift currents follow from
\eqref{in:trace-equation} and \eqref{in:bracket-shift}; the remaining
two currents follow from the metric and torsion transports.
Thus this actual-field system uses no special values of a Kerr profile.

\paragraph{Fixed regions.}
On a diamond $0\le s\le s_*$, substitute
\eqref{in:high-hodge} into \eqref{in:high-costs}. With all current
multipliers one, the resulting quadratic form is bounded by $C_m$
times their total size. The common weight
\begin{equation}\label{in:finite-weight}
 e^{-2K_{\mathrm{fin}}s},\qquad
 -L_i\log(e^{-2K_{\mathrm{fin}}s})=K_{\mathrm{fin}},
 \qquad K_{\mathrm{fin}}>C_m,
\end{equation}
absorbs the cost in both directions. Removing the weight costs only
the fixed factor $e^{2K_{\mathrm{fin}}s_*}$ on squared energies.
The recovery arguments below then give finite bounds for every
coordinate derivative. For this use extend each scalar envelope
$g_i$ by zero outside the coordinate projection of the diamond; no
spacetime field is extended. In particular, no tail assumption enters
this fixed-region argument.

\paragraph{The two limiting interactions.}
For the sharper estimate now assume the full wedge and global
summable envelopes. When $\rho_q=h_0=h_1=0$, the only unsuppressed
arrows in \eqref{in:high-costs}, after substitution of
\eqref{in:high-hodge}, are
\[
 Y_i,z_i\longrightarrow P_i,\qquad
 D\longrightarrow P_0,\qquad Y_1,z_1\longrightarrow P_\gamma.
\]
This first limiting system is not sufficient throughout
\eqref{in:high-domain}: the coordinate $y_d$ may remain bounded
while $y_l$ tends to infinity, so $h_d$ need not tend to zero there.
We will instead give the currents flowing along $d$ a large damping
on a fixed interval $-B\le y_d\le U$. Removing this part of the
weight then costs a fixed factor, independent of $T$. To identify
what the remaining currents must control, consider a second limiting
system in which the $d$-flowing currents are strongly damped.
Those currents comprise
\[
 \mathcal D_d=\{Y_l,Z_d,z_d,P_d\}
   \quad\text{and also }\{D,P_\gamma\}\text{ if }d=1.
\]
The currents flowing along $l$ comprise
\[
 \mathcal L_l=\{Y_d,Z_l,z_l,P_l\}
   \quad\text{and also }\{D,P_\gamma\}\text{ if }l=1.
\]
Delete products containing a variable of $\mathcal D_d$ and set
$\rho_q=h_l=0$, allowing $h_d$ to remain bounded. The only
remaining arrows are
\begin{equation}\label{in:asymmetric-arrows}
 z_l\longrightarrow Y_d,P_l,
 \qquad z_l\longrightarrow P_\gamma\quad\text{if }l=1.
\end{equation}
Indeed the possible term $h_dd_l^*$ either vanishes or contains
$D$ flowing along $d$; if $D$ remains in $\mathcal L_l$, its cost
vanishes. The retained $P_l$ source contains $Y_l$, which belongs
to $\mathcal D_d$, and the retained trace cost has factor $h_l$.
There is consequently no unsuppressed return to $z_l$.
Both limiting systems are triangular: their surviving interactions
can be made as small as desired by rescaling the auxiliary currents.
The actual equations retain all their terms; these two comparisons
only determine the multipliers and the directional damping.

Fix an arbitrarily small $\delta_h>0$. Keep multiplier one on each
$Y_i,Z_i$ pair and on $D^2$, multiply $z_i^2$ by $\Lambda^2$
with $\Lambda$ large, and multiply the $P$ currents by
$\epsilon_P^2$ with $\epsilon_P$ small. The resulting size is
\begin{equation}\label{in:rescaled-energy}
 \mathcal E_m=\sum_i(Y_i^2+Z_i^2)+D^2
       +\Lambda^2\sum_i z_i^2
       +\epsilon_P^2(P_\gamma^2+P_0^2+P_1^2).
\end{equation}
In the scaled variables, the first limiting comparison has coefficients
$O(\epsilon_P)$ or $O(\epsilon_P/\Lambda)$, and
\eqref{in:asymmetric-arrows} has coefficients
$O(h_d/\Lambda)$ or $O(\epsilon_P/\Lambda)$.
Choose the rescaling so that both limiting costs are at most
$\delta_h\mathcal E_m/4$ in their respective retained groups.
The principal paired angular cancellation is unchanged.
Continuity of the remaining finite quadratic forms gives a threshold
$\epsilon_*>0$ such that the full cost is small when
$\rho_q,h_0,h_1\le\epsilon_*$, and the cost involving only
$\mathcal L_l$ is small when $\rho_q,h_l\le\epsilon_*$.
This threshold is chosen before the large damping constant below.

\paragraph{Directional damping on expanding regions.}
Use the common weight
\begin{equation}\label{in:high-weight}
 \mathfrak w=\exp\{-2\delta_h(y_0+y_1)
       -2K_0(y_d+B)_+-2K_0\min(s,S_0)\}.
\end{equation}
Choose $B,S_0$ large, independently of $T$, so that for
$s>S_0,y_d<-B$, both $h_i$ and $\rho_q$ are below
$\epsilon_*$. For $s>S_0,-B\le y_d\le U$, make
$h_l,\rho_q$ small; this is possible because
$y_l=2s-y_d\ge2S_0-U$. Almost everywhere the damping in direction
$k$ is
\[
 -L_k\log\mathfrak w
 =2\delta_h+2K_0\mathbf1_{\{k=d,\ y_d>-B\}}
                         +K_0\mathbf1_{\{s<S_0\}}.
\]
Lipschitz weights suffice. With the rescaling already fixed, all
remaining interactions have bounded coefficients independent of
$K_0,T,B,S_0$. In the asymmetric late region, Cauchy--Schwarz bounds
the cross terms by a small multiple of $|\mathcal L_l|^2$ plus a
fixed large multiple of $|\mathcal D_d|^2$. Take $K_0$ large enough
to absorb the latter. On $s<S_0$ both directions receive the large
damping, and in the other late region all coefficients are small.
The order of choices is therefore
\[
 (\Lambda,\epsilon_P),\quad\epsilon_*,\quad(B,S_0),\quad K_0.
\]
Increasing $K_0$ changes none of the earlier thresholds.

The null and $s$ fluxes are bounded by the weighted initial norms.
On \eqref{in:high-domain},
$-CT\le y_d\le U$ and $-U\le y_l\le CT$, so the initial interval
has at most polynomial length. Moreover
\begin{equation}\label{in:weight-range}
 1\ge\mathfrak w\ge
 e^{-2\delta_h(U+CT)-2K_0(U+B)_+-2K_0S_0}.
\end{equation}
The last two losses are fixed. Initial higher norms are $e^{o(T)}$
by Proposition~\ref{in:entry}; since $\delta_h$ is arbitrary,
removing the weight preserves $e^{o(T)}$. In particular,
\begin{equation}\label{in:transverse-flux}
 \sup_{y_j}\|Y_i+d_i^*\|_{L^2(dy_i)}\le e^{o(T)}.
\end{equation}
For $d_0^*$ the current flows along $1$, so this is exactly the
required transverse flux in $dy_0$.

\paragraph{Recovery without dividing by the lapse.}
The estimate for $Z_i$ is not divided by the exponentially small
$\sqrt q$. Instead use the damped trace equation, with its
$C_mh_i$ homogeneous coefficient integrated along $i$. Since that
integral is uniformly bounded, integration of the damped trace equation gives
\begin{equation}\label{in:trace-convolution}
\begin{split}
 z_i(y)\le{}&e^{o(T)}e^{-\kappa(y_i+y_j)}\\
 &+C_m'\int_{-y_j}^{y_i}e^{-\kappa(y_i-y_i')}
       h_i(Y_i+d_i^*+h_iP_*)(y_i',y_j)\,dy_i'.
\end{split}
\end{equation}
The linear norm versions of the metric and torsion currents,
\eqref{in:transverse-flux}, and this convolution yield
\begin{equation}\label{in:Volterra}
 P_*(y)\le e^{o(T)}
       +C_m''\sum_i\int_{-y_j}^{y_i}
                 h_iP_*(y_i',y_j)\,dy_i'.
\end{equation}
Cauchy--Schwarz on intervals of length $O(T)$ costs only a polynomial.
There is a bounded common potential
\begin{equation}\label{in:potential}
 \Phi(y)=C_0\sum_i\int_{-\infty}^{y_i}g_i(v)\,dv
             +\frac{C_0q_0}{\kappa}(1-e^{-2\kappa s}),
 \qquad \partial_{y_i}\Phi=h_i.
\end{equation}
Multiply \eqref{in:Volterra} by $e^{-C\Phi(y)}$ and take a
supremum. Along each coordinate path,
\[
 \int_{-y_j}^{y_i}h_i(y_i',y_j)
       e^{-C(\Phi(y)-\Phi(y_i',y_j))}\,dy_i'\le C^{-1}.
\]
Choosing $C>4C_m''$ absorbs the two integrals. Hence
$P_*,z_i\le e^{o(T)}$, and \eqref{in:high-hodge} makes $R_i$
integrable along $i$ at that cost. On fixed diamonds the same argument
uses the finite flux bound from \eqref{in:finite-weight} and the
zero-extended scalar envelopes, giving finite constants.

The cross transports for $\chi_i,\ell_i+\kappa$ now give
pointwise angular bounds, initially through $H^{m-1}$.
Their homogeneous coefficient is $h_j$; the terms $qP_*$ have the
necessary spare angular derivatives, and the remaining terms
$h_i(R_j+d_j^*)+h_i h_jP_*$ integrate by the bounds just obtained.
The shift equation gives pointwise angular control through order $m$.
Since $m$ was arbitrary, every fixed angular order is available.

\paragraph{Longitudinal recovery.}
Induct on the number of longitudinal derivatives, retaining arbitrary
fixed angular orders. Proposition~\ref{in:pure-jets} estimated the first
pure derivatives at low angular order. To start the present induction,
commute their equations \eqref{in:pure-top} and \eqref{in:pure-shift}
at each fixed angular order. Their inhomogeneous terms now have the
angular bounds just proved, and their initial jets follow from the
entry tangential derivatives and the cross equations. The top transport
coefficient, including the first angular derivative of the shift, is
bounded by $C_mh_j$. Higher angular derivatives of coefficients multiply
lower angular derivatives of the pure jet. Induction in angular order therefore bounds the first
pure shape jets, then the pure lapse jets through their triangular
equation, and the pure shift jets. This starts the longitudinal
induction at every fixed angular order.

At the next level, differentiation of the cross equations gives
transports for $\mathscr L_i^k\chi_i,L_i^k\ell_i$ along $j$.
The only new pure jet on the right is the shape jet, with homogeneous
coefficient $Ch_j$; all other terms have lower longitudinal counts.
This follows from \eqref{in:pure-identities}: a derivative of either
torsion uses the torsion equation and $d\ell_i$, and curvature
variation uses $\beta_i$. Angular connection variations are expressed
by $\chi_i$ and its angular derivatives. Differentiated brackets
are angular and use only previous longitudinal orders. Derivatives of
$a$ retain its factor.

The initial pure jets follow from the noncharacteristic initial
equations, or equivalently from tangential $x$ derivatives and the
cross equations. Pointwise transport Gronwall uses only $Ch_j$,
including the first angular derivative of the shift. A subsidiary
angular induction assigns higher coefficient derivatives to already
controlled factors. Integration over length $O(T)$ of an $e^{o(T)}$
source is still subexponential. Mixed longitudinal derivatives reduce
to the cross equations and the angular commutators.

Finally recover ordinary derivatives of the shift from
\eqref{in:bracket-shift}. Pure $\partial_{y_0}^k b$ is integrated
in $y_1$; after converting Lie derivatives to coordinate derivatives,
its source contains the shift only at preceding longitudinal orders.
The other derivatives are direct, and
$a^{-1}\partial_{y_1}b=\gamma^{-1}(\xi_1-\xi_0)$ is algebraic.
The same recovery with finite bounds proves the first assertion, and
with \eqref{in:weight-range} proves the second.
\end{proof}

\subsection{Continuation and exchange of the null labels}

We now construct the full wedge. Only the fixed-region part of
Proposition~\ref{in:higher} is used for continuation. Once the wedge
exists, its global low envelopes supply the tails required for the
subexponential conclusion.

\begin{proposition}[Continuation from the full cylinder]
\label{in:continuation}
The smooth vacuum cylinder data of Proposition~\ref{in:entry} have
a smooth development on the full product wedge
$\{s\ge0\}\times\mathbb R_x\times\mathbb S^2$ in the gauge
\eqref{in:gauge}. It satisfies Propositions~\ref{in:low},
\ref{in:pure-jets}, and both conclusions of
Proposition~\ref{in:higher}. With a sufficiently small past collar,
$\mathcal S$ is a Cauchy surface for this development. The construction
does not identify this region with the whole maximal globally
hyperbolic development.
\end{proposition}

\begin{proof}
Start with smooth vacuum Cauchy evolution near $\mathcal S$ and solve
the noncharacteristic optical and angular transport equations there.
We continue by exhaustion and finite restarts.
For a target compact subdiamond work on a slightly larger region
\[
 0\le s\le s_*,\qquad y_i\le C_i-\alpha s,
\]
where $\alpha>0$ is small and the constants $C_i$ leave positive
slice length and a fixed margin around the target. The tilted sides
are strictly exiting, so all previous past-cap energies retain their
flux signs.

Bootstrap the uniform low estimates. The common finite weight
\eqref{in:finite-weight}, followed by the recovery part of
Proposition~\ref{in:higher}, bounds every fixed smooth norm up to
a putative terminal time. Since $s$ is bounded here,
$a=q_0e^{-2\kappa s}A$ has a positive lower bound; the lapse and
spatial metrics remain uniformly nondegenerate.

Restart from a slice before the alleged terminal time. Extend its
induced metric and second fundamental form smoothly across the interval
endpoints, preserving positivity and bounded requisite finite Sobolev
norms on an enlarged margin. Constraints are unnecessary outside that
margin. A reduced wave Einstein equation, initialized in wave gauge,
has a uniform short existence time from these bounds. Solve the optical
and angular transport equations with their inherited initial values.
Their transversality and invertibility on a smaller margin also persist
for a uniform short time. Higher smooth norms propagate separately.

Restrict to the exiting region. Future causal vectors obey
$dy_i\ge0$ in this gauge, so data outside the margin cannot enter
the region under consideration. Constraint and gauge propagation give
vacuum there. Geometric uniqueness, followed by uniqueness of the
optical functions and angular transport, identifies the restarted and
previous solutions. Slightly larger diamonds supply the endpoint jets
and margins for each restart, patchwise on the sphere. Thus no finite
terminal time occurs on a compact target diamond.

Exhaustion gives the full smooth product wedge. Increasing-domain
suprema and their cell sums preserve the uniform low bounds and define
global summable envelopes $g_i$. Only now use their tails in the
$B,S_0$ choices of \eqref{in:high-weight}; this proves the expanding
subexponential bounds without circularity.

Finally $s$ is temporal. Causality in \eqref{in:gauge} implies
$|dx/ds|\le1$, and angular speeds are bounded on compact $(s,x)$
ranges. A past causal curve from the wedge therefore reaches
$\mathcal S$ or has an ordinary limit where it continues; it cannot
escape at finite decreasing $s$. Choose a local past collar of
possibly varying smooth thickness in $x$, with sufficiently small
Lipschitz slope within the local existence region. Future causal
curves in that collar reach the cylinder before leaving it. The
cylinder is consequently Cauchy for this development, which can be
attached along all of $\mathcal S$ by geometric Cauchy uniqueness.
\end{proof}

\begin{lemma}[The same wedge with exchanged labels]\label{in:label-swap}
The two optical labels may be exchanged while retaining the same
initial cylinder and the same product spacetime. In the corresponding
one-shift angular gauge, all conclusions of
Propositions~\ref{in:entry}--\ref{in:continuation} remain valid,
with the same common $\kappa$ and with a frozen physical profile
depending only on the new $y_0$.
\end{lemma}

\begin{proof}
Negate the prepared $x$ on $\mathcal S$, including its bounded joins.
The old optical functions exchanged have the required eikonal initial
data, so optical uniqueness identifies them with the new pair.
For the angular labels use the exact preparation that eliminates the
new second shift. In the existing geometry, transport these labels
along the old $L_0$. At every finite point this is a finite-duration
smooth flow on a compact sphere over a finite coordinate interval,
hence a sphere diffeomorphism with smooth inverse. It defines the new
angular chart on the same full wedge before any uniform estimate on
that flow is assumed.

To verify the initial comparison, reverse $T_K,\varphi$
simultaneously at both ends, keep $Q$, and use $-I,-I_{\mathrm{rot}}$
with the new transport map $H$ and the corresponding bounded angular
joins. Thus the surface $s=0$ is unchanged and its prepared $x$ is
exactly negated. Use the reversed $I$ to define a single physical
parameter profile along the new $y_0$. The exact profile estimates
and the trace proof of Proposition~\ref{in:entry} apply with this
other frozen angular map. Apply the finite-diamond a priori estimates
to the new fields and exhaust. This reestablishes all uniform low and
subexponential high bounds on the same spacetime, without requiring
an a priori uniform estimate for the long angular coordinate change.
\end{proof}
